\chapter{Three-Manifolds with Positive Ricci Curvature, and Beyond}

\section{Three-manifolds with positive Ricci curvature, and beyond}

\subsection{Hamilton's theorem}

\begin{theorem}
\label{topping-ch10-hamilton-theorem}
\source{topping:page-0124,topping:page-0125,topping:page-0126}
\uses{topping-maxpr-finite-time-blowup, topping-ch8-blow-up-singularities, topping-ch9-roundness-estimate, topping-ricci-flow-uniqueness}
\dcref{ch10:1.1}
\group{positive-ricci-three-manifolds}

Let $(\mathcal{M},g_0)$ be a closed three-dimensional Riemannian manifold with
positive Ricci curvature, and let $g(t)$ be its Ricci flow with $g(0)=g_0$ on
the maximal interval $[0,T)$. There exist a metric $g_\infty$ on $\mathcal{M}$
of constant positive sectional curvature, times $t_i\uparrow T$, points
$p_i\in\mathcal{M}$, a point $p_\infty\in\mathcal{M}$, and a constant $c>0$
such that, for $t\leq 0$, the rescaled flows
\[
g_i(t)=\abs{\Rm(p_i,t_i)}\,g\!\left(t_i+
\frac{t}{\abs{\Rm(p_i,t_i)}}\right)
\]
satisfy
\[
(\mathcal{M},g_i(t),p_i)\longrightarrow
(\mathcal{N},(c-t)g_\infty,p_\infty)
\]
on $(-\infty,0]$.
\end{theorem}

\begin{proof}
  \uses{topping-maxpr-finite-time-blowup, topping-ch8-blow-up-singularities, topping-ch9-roundness-estimate, topping-ricci-flow-uniqueness}
\sketch Compactness and positive Ricci curvature give constants
$0<\beta\leq B<\infty$ with
$\beta g_0\leq\Ric(g_0)\leq Bg_0$. Hence $R(0)\geq3\beta>0$, so the
maximal time $T$ is finite. The blow-up compactness theorem supplies the
sequences $t_i\uparrow T$, $p_i$, the rescaled flows $g_i(t)$, and a limit
$(\mathcal{N},\hat g(t),p_\infty)$ for $t\in(-\infty,b)$.

The roundness estimate gives, for every $\gamma>0$,
\[
\left|\Ric-\frac13Rg\right|\leq\gamma R+
\frac{M(\beta,B,\gamma)}{\abs{\Rm(p_i,t_i)}}
\]
on the rescaled flows at time zero. The error tends to zero, so the limit
satisfies $|\Ric-\frac13R\hat g|\leq\gamma R$ for every $\gamma>0$ and is
Einstein. In dimension three it therefore has constant sectional curvature.
The preserved nonnegative sectional curvature and the normalization
$\abs{\Rm(\hat g(0))(p_\infty)}=1$ make that curvature strictly positive.

Myers' theorem makes $\mathcal{N}$ compact, so the limit is closed and the
flow convergence identifies it with $\mathcal{M}$. Backward uniqueness gives
$\hat g(t)=(c-t)g_\infty$ for a positive multiple $g_\infty$ of $\hat g(0)$,
which is the asserted convergence.
\end{proof}

\begin{corollary}
\label{topping-ch10-space-form-corollary}
\source{topping:page-0126}
\uses{topping-ch10-hamilton-theorem}
\dcref{ch10:1.2}
\group{positive-ricci-three-manifolds}

Every closed Riemannian three-manifold with positive Ricci curvature admits a
metric of constant positive sectional curvature. If it is simply connected,
then it is $S^3$.
\end{corollary}

\subsection{Beyond the case of positive Ricci curvature}

\begin{theorem}
\label{topping-ch10-hamilton-ivey-theorem}
\source{topping:page-0126}
\uses{topping-ch9-ode-pde-theorem}
\dcref{ch10:2.1}
\group{positive-ricci-three-manifolds}

For every finite $B$ and every $\varepsilon>0$, there is
$M=M(\varepsilon,B)<\infty$ such that, whenever
$(\mathcal{M},g_0)$ is a closed three-manifold with
$\abs{\Rm(g_0)}\leq B$, its Ricci flow $g(t)$ with $g(0)=g_0$ satisfies
\[
E+(\varepsilon R+M)g\geq0,
\]
where $E$ is the Einstein tensor.
\end{theorem}

In dimension three, this estimate implies
\[
\lambda_1\geq-\varepsilon R-M,
\qquad
R=2(\lambda_1+\lambda_2+\lambda_3).
\]
Thus a substantially negative lowest sectional curvature requires the largest
sectional curvature to be correspondingly larger.

\begin{corollary}
\label{topping-ch10-blowup-weakly-positive}
\source{topping:page-0127}
\uses{topping-ch10-hamilton-ivey-theorem}
\dcref{ch10:2.2}
\group{positive-ricci-three-manifolds}

For a Ricci flow on a closed three-dimensional manifold, every blow-up limit
\[
(\mathcal{M},g_i(t),p_i)\longrightarrow(\mathcal{N},\hat g(t),p_\infty)
\]
obtained as in Theorem~\ref{topping-ch8-blow-up-singularities} has weakly positive sectional curvature on
$(\mathcal{N},\hat g(t))$ for all $t\in(-\infty,b)$.
\end{corollary}
